\documentclass[11pt,a4paper]{article}
\usepackage[T1]{fontenc}\usepackage{lmodern,amsmath,amssymb,amsthm,mathtools,microtype}
\usepackage[margin=27mm]{geometry}\usepackage[colorlinks=true,linkcolor=blue,urlcolor=blue,citecolor=blue]{hyperref}
\newtheorem{theorem}{Theorem}[section]\newtheorem{lemma}[theorem]{Lemma}\newtheorem{corollary}[theorem]{Corollary}\newtheorem{proposition}[theorem]{Proposition}
\newcommand{\Q}{\mathbb Q}\newcommand{\Z}{\mathbb Z}\DeclareMathOperator{\Tr}{Tr}
\setlength{\emergencystretch}{2em}
\title{Trace heights on quadratic twists with $j=1728$}
\author{}\date{13 November 2026}
\begin{document}\maketitle
\begin{abstract}
For a quadratic number field $K$ and $d\in K^\times$ with
$d\ne\sigma(d)$ and $d\ne-\sigma(d)$, the rational trace of the
abscissa on $y^2=x^3-d^2x$ controls the canonical height of every
$K$-rational point with finite abscissa. The additive constant is
ineffective. The proof excludes all positive-dimensional abelian
translates in the relevant finite fibers and checks the projective
embedding required by von Buhren's family height theorem. No rank-one
hypothesis or square-abscissa condition is needed for this height result.
For the curve used in the Diophantine construction, the conclusion
therefore covers half-integral indices and torsion translates. We also
supply a reciprocal-coordinate description of an ordered pair of points
and its difference, including the identity and tangent cases.
\end{abstract}

\section{The height theorem}
Heights are absolute logarithmic Weil heights. For $a/b\in\Q$ reduced,
$b>0$, write $h(a/b)=\log\max(|a|,b)$. Canonical heights on elliptic
curves use the standard symmetric divisor at the identity; changing
its normalization only changes the positive constants below.

\begin{theorem}\label{thm:height}
Let $K/\Q$ be quadratic with nontrivial automorphism $\sigma$. Let
$d\in K^\times$ satisfy $d\ne\sigma(d)$ and $d\ne-\sigma(d)$, and set
$E_d:y^2=x^3-d^2x$. There are constants $A>0$ and $J\ge0$, depending
only on $K,d$, such that every $Q\in E_d(K)\setminus\{O\}$ satisfies
\[
 \widehat h(Q)\le A h\bigl(\Tr_{K/\Q}x(Q)\bigr)+J.
\]
The theorem is unconditional. The argument does not give an effective
numerical value of $J$.
\end{theorem}

\begin{lemma}\label{lem:fibers}
On $E_0^2$, where $E_0:Y^2=X^3-X$, let
$F(Q_1,Q_2)=a x(Q_1)+b x(Q_2)$ with $ab\ne0$ and $a\ne\pm b$.
No translate of a positive-dimensional abelian subvariety which meets
$(E_0\setminus\{O\})^2$ has constant restriction of $F$.
\end{lemma}
\begin{proof}
The whole surface does not have constant $F$. A proper
positive-dimensional abelian subvariety is an elliptic curve. If one
projection of its translate is constant, the other projection is
nonconstant and $F$ is not constant. Otherwise, after an isogeny from
$E_0$ to this elliptic curve, the two projections have the form
$\alpha z+q_1$ and $\beta z+q_2$, with nonzero endomorphisms
$\alpha,\beta$ of $E_0$.

If $F$ were constant, its two summands would have identical pole
supports. Their poles are all double, since the isogenies are etale
in characteristic zero. Choose a common pole and translate the
parameter so it is zero. The projections become $\alpha z,\beta z$;
the equality of pole supports gives $\ker\alpha=\ker\beta$. Factoring
through their common quotient gives $\beta=u\alpha$ for an
origin-preserving automorphism $u$ of $E_0$. On this short
Weierstrass model such an automorphism sends $x$ to $x$ or $-x$.
For example, its action on the linear system $|2O|$ is affine in $x$;
preservation of the three finite branch points $-1,0,1$ forces its
translation to be zero and its slope to be $1$ or $-1$.
Consequently the restriction is $(a+b)x(\alpha z)$ or
$(a-b)x(\alpha z)$, a nonconstant function under the hypotheses.
\end{proof}

\begin{proof}[Proof of Theorem~\ref{thm:height}]
Choose square roots of $d$ and $d'=\sigma(d)$ in a fixed number-field
extension. The pointed isomorphisms to $E_0$ send $(x,y)$ to
$(x/d,y/d^{3/2})$ and $(x/d',y/(d')^{3/2})$, respectively. The image
of a point and its conjugate is therefore a point of $E_0^2$ at which
\[
                 F=d x_1+d'x_2=\Tr_{K/\Q}x(Q).
\]
Take the integral reduced projective closure $\mathcal V$ of the
graph of $F$ in $E_0^2\times\mathbf P^1$. Its dimension is two.
The set of images of $E_d(K)$ is a finitely generated subgroup
$\Gamma$ of $E_0^2(\overline\Q)$, by Mordell--Weil and the fact that
the two pointed maps commute with the group operations.

We apply the published Theorem~1.1 of von Buhren~\cite{VB}. It bounds
canonical height linearly in the parameter height, with a fixed
additive constant for each finite-rank subgroup, away from the union
of positive-dimensional abelian translates contained in the fiber.
A point of the graph above an actual finite value of $F$ lies outside
that union: a translate through it would meet the domain of $F$ in a
dense open subset and make $F$ constant there, contrary to
Lemma~\ref{lem:fibers}. This also treats possible components of the
graph over the indeterminacy locus; none changes the argument through
our point in the domain.

Here are the remaining embedding and subgroup checks. Embed $E_0$ as
its smooth plane cubic by $\mathcal O(3O)$, and embed its square by
Segre in $\mathbf P^8$. For $n\ge1$ the sequence
\[
0\longrightarrow\mathcal O_{\mathbf P^2}(n-3)
 \longrightarrow\mathcal O_{\mathbf P^2}(n)
 \longrightarrow\mathcal O_{E_0}(3nO)\longrightarrow0
\]
and $H^1(\mathbf P^2,\mathcal O(n-3))=0$ prove surjectivity of the
restriction of sections. Tensoring the two surjections and using
Segre proves condition (1.1) in von Buhren's statement for every
$n\ge1$. The line bundle on the surface is ample and symmetric.
On the base use $\mathcal O_{\mathbf P^1}(1)$ with its ordinary height.
Every point of $\Gamma$ lies in the source set $C(\Gamma,\epsilon)$
by taking its error point to be zero. Thus no uniformity in a varying
group or field is being invoked.

The canonical height for the chosen product embedding is a fixed
positive multiple of the sum of the two elliptic canonical heights.
Pointed isomorphism invariance and Galois invariance identify these
with the canonical height of $Q$. The source theorem now gives the
displayed inequality. Its additive constant depends on this fixed
$\Gamma$ and is not made effective by the source. Increase constants
if necessary to make $A>0$ and $J\ge0$.
\end{proof}

\begin{corollary}\label{cor:index}
Let $K=\Q(\sqrt2)$, $d=3-\sqrt2$, and $P=2R_0$, where
$R_0=(1+\sqrt2,2\sqrt2)$. Retain the arithmetic premise
$2E_d(K)=\Z P$ with $P$ non-torsion. Every $Q\in E_d(K)$ has a unique
index $n\in\tfrac12\Z$, defined by $Q=[2n]R_0+T_0$ with $2T_0=O$.
There are fixed $A>0,J\ge0$ such that
\[
 n^2\le A h\bigl(\Tr x(Q)\bigr)+J\qquad(Q\ne O).
\]
This includes finite torsion points, with index zero.
\end{corollary}
\begin{proof}
If $2Q=mP$, then $2(Q-mR_0)=O$. Uniqueness follows from
non-torsion of $R_0$. The canonical height of $Q$ equals
$n^2\widehat h(P)$. Apply the theorem and divide by its positive
constant $\widehat h(P)$; $d+d'=6$ and $d-d'=-2\sqrt2$ are nonzero.
\end{proof}

\section{Two reciprocal points and their difference}
This section is algebraic and does not assume a group-rank statement.
Work over any field of characteristic zero, with $D\ne0$, on
$E:y^2=x^3-Dx$. The rational functions
\[
          f=1/x,
          \qquad h=-y/x^2
\]
extend by $(f,h)=(0,0)$ at $O$. They identify
$E\setminus\{(0,0)\}$ with $h^2=f-Df^3$. Indeed the affine map
$(f,h)\mapsto(-Df,-Dh)$ gives translation by $(0,0)$ on $E$.

\begin{proposition}\label{prop:pair}
The three equations
\begin{align}
 h^2&=f-Df^3,\\
 s(f-g)&=h+k,\\
 s(h-k)&=1-D(f^2+fg+g^2)\label{eq:tangent}
\end{align}
describe ordered pairs $U,V\in E\setminus\{(0,0)\}$ with $U\ne V$,
and their unique secant or tangent parameter $s$. Their difference
is the finite point
\[
 A=U-V=(c,y_A),\qquad
 c=s^2+D(f+g),\quad y_A=Dh-s(Df+c).
\]
Either $U$ or $V$ may be $O$. Finite torsion differences are permitted.
\end{proposition}
\begin{proof}
Put $r_1=s(f-g)-h-k$, $r_2=s(h-k)-1+D(f^2+fg+g^2)$ and
$r_3=h^2-f+Df^3$. Exact polynomial identities give
\begin{align*}
 k^2-g+Dg^3&=r_3-(f-g)r_2+(h-k)r_1,\\
 y_A^2-c^3+Dc&=D^2r_3-D(Df+c)(r_2-sr_1).
\end{align*}
Thus both reciprocal points and the asserted finite point are on the
curve. If $f\ne g$, translate both points by $(0,0)$. The line
joining the first translated point to the negative of the second has
slope $s$; the chord law gives exactly $(c,y_A)$. Translation by this
two-torsion point does not alter their difference.

If $f=g$, the equations give $k=-h$. If also $h=0$, the cubic gives
$f=0$ or $Df^2=1$, whereas~\eqref{eq:tangent} would give
$1-3Df^2=0$, a contradiction in either case. Hence $h\ne0$ and
$2sh=1-3Df^2$. This is the tangent slope at the first translated
point, so the tangent law gives the same formula. In particular
$U\ne V$ in every case. Conversely these secant and tangent
calculations supply $s$ for every allowed ordered pair. The identity
case is the ordinary affine point $(f,h)=(0,0)$ in this translated
model and is already covered.
\end{proof}

\section{Ratio equations and formal precision}
If $H$ is finite with $S=y(H)/x(H)\ne0$, its formal parameter is
$z(H)=-1/S$. The two equations
\[
              hL+fS=0,\qquad (1-Df^2)L+hS=0
\]
give $L=z(U)/z(H)$ whenever $U$ is non-torsion or $O$; at $O$ the
second equation gives $L=0$. At a nonidentity formal point $h\ne0$,
so the first equation gives the ratio and the second follows from the
cubic. We impose both equations in the quartic construction and do
not divide by $h$ on other components.

For completeness, at a good unramified prime $p>5$, $D$ a unit, the
invariant differential and formal logarithm in $z=-x/y$ are
\[
 \omega=\frac{dz}{\sqrt{1+4Dz^4}},\qquad
 \log_E(z)=\sum_{j\ge0}\frac{(-1)^j\binom{2j}{j}D^j}{4j+1}z^{4j+1}.
\]
The first identity follows from
$x=(1+\sqrt{1+4Dz^4})/(2z^2)$. The binomial coefficients are integers.
For $j\ge1$, $v_p(4j+1)\le4j-4$, so if $v_p(z)\ge1$ all omitted
terms have relative valuation at least $4v_p(z)$. Consequently
$\log_E(z)=z(1+O(p^{4v_p(z)}))$. At $O$ both logarithm and parameter
are zero; the multiplicative unit is assigned the value one.
This supplies the order-four relative precision used by the new
quartic without imposing a positive lower bound on a numerator index.

\section{Verification and source scope}
The exact code accompanying the 81-variable quartic checks the two
polynomial identities for the pair symbolically. A separate
standard-library program checks 702 ordered pairs over $\Q(\sqrt2)$,
including 24 tangent pairs, 52 pairs involving the identity, 78
nonzero torsion differences, and 352 half-integral index differences.
It also checks 1300 regular ratio instances. These are finite examples,
not a substitute for the universal proof or for Mordell--Weil.

The final published statement of von Buhren's Theorem 1.1 and its
embedding condition were checked in text and on the image of printed
page 308. The application above is proved here; the internal proof of
the published theorem is not represented as independently re-audited.
No $abc$, prime-pattern or pole-realization hypothesis enters
Theorem~\ref{thm:height}; the rank premise enters only its index
corollary.

\begin{thebibliography}{9}
\bibitem{VB} J\'er\^ome von Buhren, \emph{Borne de hauteur semi-effective
pour le probl\`eme de Mordell--Lang dans une vari\'et\'e ab\'elienne},
Journal de Th\'eorie des Nombres de Bordeaux 29 (2017), 307--320,
Theorem 1.1, p.~308. DOI: 10.5802/jtnb.981.
\url{https://www.numdam.org/article/JTNB_2017__29_1_307_0.pdf}.
\end{thebibliography}
\end{document}
